% !TEX root = ../../hw02.tex
% Homework 02, Exercise 6. Shared problem statement.
\begingroup
\renewcommand{\labelenumi}{(\alph{enumi})}
\begin{exercise}
  \label{ex:hw02:newton-iteration}
  \solutionlink{hw02:newton-iteration}
  In class, we discussed the iteration process
  \[
    x_{n+1}=a+\frac{1}{x_n},\quad n=0,1,2,\ldots,
    \quad a\geq 1,\quad x_0>0,
  \]
  whose limit gives us the continuous fraction
  \[
    b=a+\frac{1}{a+\frac{1}{a+\frac{1}{a+\cdots}}}.
  \]
  That is, $b$ is a fixed point of the function
  \[
    f(x)=a+\frac{1}{x},\quad x>0,
  \]
  which is realized as a root of the function $F(x)=x-f(x)$ $(x>0)$.
  \begin{enumerate}
    \item Show that the iterative algorithm
      \[
        x_{n+1}=\frac{ax_n^2+2x_n}{1+x_n^2},
        \quad n=0,1,2,\ldots,\quad x_0>0,
      \]
      arises from Newton's method that may be used to approach $b$ of our problem.
    \item Assume that $\{x_n\}$ in (a) converges.
      Verify that the limit is indeed the continuous fraction $b$ defined above.
    \item Take $a=\frac{3}{2}$. Prove $0<x_n\leq 2$ for $n=1,2,\ldots$.
    \item Prove that $\{x_n\}_{n=1,2,\ldots}$ increases for any $x_0\ne 2$
      and is constant when $x_0=2$.
    \item Show that $\{x_n\}$ converges to $2$ with any initial state $x_0$
      (thereby establishing the ``global convergence'' of the algorithm).
  \end{enumerate}
\end{exercise}
\endgroup
